\section{Vulnerability Identification}
\label{detection}

The task of identifying inputs that trigger sanitizers is carried out mostly using dynamic analysis techniques.
This is largely dictated by the structure of the competition, since the targets are provided in OSS-Fuzz format with specific harnesses to access the application functionality.

At a high level, there are a number of fuzzers that operate in different ways: The \afluzzer agent uses a custom version of the AFL++ tool~\cite{aflplusplus} to fuzz a C/C++ application; The \jazzmine agent uses a custom version of the Jazzer tool~\cite{jazzer} to fuzz a Java application.

In addition, there is a grammar-focused agent, called \grammarguy, whose task is to develop grammars that can be used to generate input seeds that conform to specific formats.

Finally, there is an agent, called \quickseed, who uses LLM-based reasoning to produce interesting inputs for a harness.

Every agent that can produce interesting seeds (i.e., both the fuzzers and the input generators) contributes their findings to a seed repository called \seedsync, which takes care of distributing the seed across all the fuzzing agents.

However, not all vulnerability finding is performed using fuzzing.
There is an agent, called \lluminar, that looks at every function using a custom LLM, and determines if the function is vulnerable.
If that's the case, then the LLM reasons about how to provide an input that can reach the vulnerable function using LLM-based reasoning. 

%ijon system will be included 

\subsection{Coverage-Guided Fuzzing}
As a baseline, we use \afl~\cite{aflplusplus} and \libfuzzer~\cite{libfuzzer} to fuzz C and C++ codebases, and \jazzmine{}, our Java fuzzer based on \jazzer~\cite{jazzer}, for Java projects..
Instances of \afl\ and \jazzmine\ are stochastically configured from a predefined configuration space inherent to each tool to ensure diversity of fuzzing instances.
Examples of such configurations include fuzzing timeouts, CmpLog\cite{aschermann2019redqueen}, string dictionary usage, whether to use grammar-fuzzing, etc.
This diverse set of fuzzing configurations improves both the bug-finding ability of our system and the resilience of the system to unknown targets.
All instances of a given fuzzer share discovered seeds between each other, both within a given fuzzing node as well as across fuzzing nodes.
\subsection{AFluzzer}
\label{afluzzer}

As a baseline, we use \afl~\cite{aflplusplus} and \libfuzzer~\cite{libfuzzer} to fuzz C and C++ codebases, and \jazzmine{}, our Java fuzzer based on \jazzer~\cite{jazzer}.
Each instance of \afl\ and \jazzmine\ selects a random configuration from a predefined configuration space inherent to each tool.
Examples of such configurations include fuzzing timeouts, CmpLog\cite{aschermann2019redqueen}, string dictionary usage, whether to use grammar-fuzzing, etc.
This diverse set of fuzzing configurations improves both the bug-finding ability of our system and the resilience of the system to unknown targets.
All instances also share seeds amongst each other.
\subsection{\jazzmine}

\jazzmine is the agent responsible for fuzzing Java targets using a modified version of the Jazzer tool~\cite{jazzer}.

As the first step, \jazzmine uses CodeQL to extract all of the interesting strings from the target, in order to build a dictionary, which will be used during the fuzzing process.

Then, \jazzmine retrieves from the \builder a version of the target with all the sanitizers provided natively by Jazzer and, in addition, extended custom-built sanitizers.
These sanitizers are ``loosened'' version of the original sanitizers that allow for less strict match between the values in the inputs and the values used in monitored instructions of the application.
Which sanitizers are used in a specific execution can be selected through an environment variable (choice of the original set, the extended one, or both).

% \TBD: Explain why the sanitizers have been loosened with an example.
Unlike sanitizers developed for compiled languages, Java sanitizers are highly bug class-specific and employ a fundamentally different approach. These sanitizers operate by hooking sink functions and performing validation checks on data as it reaches these critical points as function arguments. However, their implementation relies on hardcoded strings known as sentinels for data control validation. A prime example is the working of OS command injection sanitizer, which hooks the \textbf{start} function of the \textbf{java.lang.ProcessImpl} class in the JDK and subsequently verifies whether this function can be invoked with \textbf{"jazze"} as the command argument. While this approach proves effective in scenarios with straightforward Input-to-State Correspondence, it becomes significantly more challenging when compression, encoding, or other obfuscation techniques are introduced. Under these conditions, the fuzzer struggles to accurately determine and generate the precise input required to trigger the sanitizer checks, since it relies solely on fuzzer mutation and execution to build the Input-to-State Correspondence. 

To address this limitation, we introduce an LLM agent that figures out the Input-to-State Correspondence and generates the required input to reach the state with the correct sentinel.  We modify the sanitizer to trigger when the hooked function is called with fuzzer artifacts, which are strings or characters that have a high probability of being generated by the fuzzer's mutator, thereby expanding or "loosening" the conditions under which the sanitizer crashes.

Then, \jazzmine calls \grammarguy and determines if a grammar has been produced for the relevant harness.

%\TBD: What is actually done with this grammar?

This grammar generates test inputs that are validated against the original sanitizer; this is the step which is used to filter out all the false positive crashes that might have been introduced due to the loosening of the sanitizer. We use the original sanitizer's crash as the ground truth for the LLM agent, if it is able to crash, then we mark the seed as the crashing input and send it to POV guy. If the agent fails to produce the crashing input within the allocated budget or iterations, we preserve the grammar and send it to \revolver for fuzzing.

There are three possible campaign modes:

\begin{enumerate}
\item original mode;
\item with extended sanitizers;
\item with grammar-based mutators (if a grammar was generated for the harness).
\end{enumerate}

The seeds for each harness are synchronized across campaigns. 
In addition, \jazzmine minimizes the seed corpus for all the fuzzers running on a shared Kubernetes node.

\jazzmine also sync seeds across nodes, and receives seeds from \grammarguy, \quickseed, and \discoveryguy.

%We also have the "secret" seeds from OSS-Fuzz.

When a crash is found, the harness and seed is sent to \povguy.

%When a crash is based on an extended sanitizer, it is sent to \grammaragentlowsanreproducer, who is responsible for making the same crashing input trigger the corresponding standard sanitizer.














\subsection{Grammar-based Fuzzing}

\subsubsection{NAUTILUS integration and extensions}

A large tenet of ARTIPHISHELL's bug-discovery pipeline is centered around grammar-based fuzzing. For ARTIPHISHELL we opted to build our tooling on top of the Nautilus-fuzzer. Nautilus supports two different grammar formats: a python representation and a JSON-based one. While the JSON-based format is more widely supported, e.g. by AFL++, atnwalk and others, we chose to build our tooling around the python representation of the grammar. The python format is advantageous in terms of the expressiveness of the grammars as well as the increased familiarity of Large Language Models with python.

An example of both the JSON and a python grammar can be compared in Figure~\ref{fig:grammar-format-comparison}.
%
\begin{figure*}
\centering
\begin{tabular}{@{}p{0.48\textwidth}@{\hskip 15pt}p{0.48\textwidth}@{}}
%
\textbf{Nautilus JSON Format}
\begin{minted}[fontsize=\footnotesize, baselinestretch=1.0, frame=single, framesep=3mm]{json}
{ "rules": {
    "START": [[{"NT": "HTML"}]],
    "HTML": [
      [ {"T": "<html>"},
        {"NT": "HEAD"},
        {"NT": "BODY"},
        {"T": "</html>"}
      ]],
    "HEAD": [
      [{"T": "<head>"},
        {"NT": "TITLE"},
        {"T": "</head>"}
      ]],
    "TITLE": [
      [{"T": "<title>"},
        {"NT": "TEXT"},
        {"T": "</title>"}
      ]],
    "BODY": [
      [{"T": "<body>"},
        {"NT": "CONTENT"},
        {"T": "</body>"}
      ]],
    "CONTENT": [
      [{"NT": "PARAGRAPH"}],
      [{"NT": "HEADING"}],
      [{"NT": "PARAGRAPH"}, {"NT": "CONTENT"}]],
    "HEADING": [
      [{"T": "<h1>"},
        {"NT": "TEXT"},
        {"T": "</h1>"}
      ]],
    "PARAGRAPH": [
      [{"T": "<p>"},
        {"NT": "TEXT"},
        {"T": "</p>"}
      ]],
    "TEXT": [
      [{"T": "Hello World"}]
    ]}}
\end{minted}
%
&
\textbf{Nautilus Python Format}
\begin{minted}[fontsize=\footnotesize, baselinestretch=1.0, frame=single, framesep=3mm]{python}
# Basic HTML Grammar
ctx.rule("START", b"{HTML}")

# HTML structure
ctx.rule("HTML", b"<html>{HEAD}{BODY}</html>")

# Head section
ctx.rule("HEAD", b"<head>{TITLE}</head>")
ctx.rule("TITLE", b"<title>{TEXT}</title>")

# Body section
ctx.rule("BODY", b"<body>{CONTENT}</body>")

# Content alternatives
ctx.rule("CONTENT", b"{PARAGRAPH}")
ctx.rule("CONTENT", b"{HEADING}")
ctx.rule("CONTENT", b"{PARAGRAPH}{CONTENT}")

# Elements
ctx.rule("HEADING", b"<h1>{TEXT}</h1>")
ctx.rule("PARAGRAPH", b"<p>{TEXT}</p>")

# Text content
ctx.literal("TEXT", b"Hello World")
\end{minted}
%
\end{tabular}
\caption{Side-by-side comparison of Nautilus supported JSON and Python format for grammar representation}
\label{fig:grammar-format-comparison}
\end{figure*}
%
First, the python representation supports the definition of arbitray "producer"-functions in python. This allows the LLMs to both leverage their familiarity of writing python code as well as to bring to bear the full expressive power of the python standard library.
This enables us to easily represent complex encodings, e.g., compression, base64, tar-archives and many more such common representations natively.
These encoder functions allow the LLMs to understand and represent multi-layered file format encodings and build them into the grammars they write, while focusing the generative power of the grammar on the inner fields the application is finally parsing.

We show an example grammar for a bug in the AIxCC organizer's Tika exemplar that makes use of this functionality to produce valid HTML pages with fully valid embedded and base64-encoded PNG images. (\TBD: fit the grammar from comment to text / appendix)
%
% Put the thing in the appendix. I don't know how else to make it look somehwhat acceptable
%
The version of Nautilus in ARTIPHISHELL was heavily customized and extended with the following additions:

\paragraph{ANYRULE} rules allow the fuzzer to attempt to insert derivations from any rule in the grammar in place of the "correct" derivations according to the grammar structure.
This extension breaks NAUTILUS' strict adherence to the grammar derivation tree in order to maximally use the coverage-guidance.
Real-world file formats oftentimes reuse encodings, inner file structures and data sequences in multiple contexts.
If the fuzzer finds that such a case leads to higher coverage, e.g., in an incomplete grammar written by an LLM, then ANYRULE-rules allow the fuzzer to further explore this embedded functionality everywhere in the target format.

\paragraph{IMPORT} rules allow us to easily represent recursively nested file formats as they are commonly found in real-world file formats.
Examples of this are, e.g. Microsoft Powerpoint documents which are zip-files containing multiple embedded XML documents and HTML documents that can directly embed various image file formats (PNG, JPEG, WEBP, etc.) and scripting language snippets.

\paragraph{Binary encodings of common integer types} Binary file formats often contain packed integer fields.
We have added native support for these encodings and added mutations that specifically operate on these fields to exploit this knowledge, e.g. increment, decrement, etc.

\paragraph{Support for byte-mutations} Nautilus supports variable fields mainly via their support for regular-expression rules that describe the structure of the value being synthesized.
These fields are not mutated, but instead generally regenerated from scratch each time, which prevents the fuzzer from incrementally discover internal structures.
Our version of nautilus additionally supports arbitrary BYTES fields of given lengths and during mutation it uses traditional fuzzing mutations on the previous value instead of regenerating the value from scratch.
This allows our NAUTILUS to discover internal file formats through repeated coverage-guided mutations.

\subsubsection{Grammar-fuzzing in ARTIPHISHELL}

The main components that perform and facilitate grammar-fuzzing in ARTIPHISHELL are:

\newpar{Grammar[Guy/Agent/Roomba]} LLM agents that automatically write and augment grammars to cover as much code in the target application as possible.

\newpar{Grammar-composer} Detects known file formats and encodings that are represented in a given full grammars as well as its internal sub-rules.
If known file formats are discovered as likely matches, it will produce new grammars by splicing grammars from  our curated corpus of pre-written grammars into these grammars.
This allows us to detect what format an LLM was attempting to represent in its rules and immediately expand the potential coverage reached by that grammar to include the full file format.
Furthermore, it will also attempt to splice in custom "corpus-grammars" that allow our grammar-fuzzers to make use of our pre-loaded fuzzing corpora inside nested grammar structures.

\newpar{\revolver} Custom mutator library used by our fuzzers to automatically synthesize and mutate seeds from multiple grammars at once.

\newpar{Fuzzing: AFL and Jazzer} We have added configurations to our AFL and Jazzer configuration selections that incorporate grammar fuzzing. Any new fuzzing instance will have a certain probability of being selected as a grammar-fuzzing instance during launch (33\%).


\newpar{\seedsync} Our seed syncing agent maintains both traditional fuzzing corpora as well as the synchronization of our grammar-fuzzing instances. This includes both the serialized derivation trees (RONs) as well as any newly discovered grammars.


\subsubsection{Grammar Composer}

While single-grammar fuzzing effectively explores individual format spaces, real-world vulnerabilities often hide at format boundaries where parsers handle nested data structures.
Traditional fuzzers miss these bugs because they cannot adapt when encountering embedded formats -- such as Base64-encoded images within HTML or EXIF metadata within JPEGs.
Our grammar composition system addresses this limitation through runtime detection and composition of grammar rules based on observed data patterns.

The grammar composer operates by monitoring the fuzzer's output queue and identifying when grammar rules generate data resembling other known formats.
It maintains a \textbf{fingerprint database} that tracks patterns in the data each grammar rule produces, using 64-bit bloom filters to efficiently detect recurring byte sequences.
When a rule consistently generates data matching a known format (e.g., JPEG magic bytes appearing in Base64 output), the composer records this as a \textbf{composition opportunity}.

Upon detecting such opportunities, the composer creates enhanced test cases by splicing the appropriate format grammar into the derivation tree.
For instance, when HTML grammar rules produce Base64-encoded image data, the composer: (1) identifies the composition opportunity, (2) replaces the rule output with the full derivation tree (Base64-encoded) for the image, (3) updates the grammar rules accordingly, and (4) return the spliced test case to the fuzzer queue for further mutation.
This process preserves the fuzzer's evolved test structure while exposing previously unreachable code paths.

The system achieves this through four operations:

\begin{itemize}
\item \textit{Fingerprint collection}: Extracts 4-byte prefixes from rule outputs and updates bloom filters.

\item \textit{Pattern detection}: Identifies format types via MIME detection (libmagic) or fingerprint matching against the reference fingerprint database.

\item \textit{Tree splicing}: Replaces rule outputs with full derivation trees from the detected grammar, handling encoding transformations.

\item \textit{Coverage filtering}: Executes both original and spliced variants, keeping only those that discover new code paths.
\end{itemize}

This pattern-recognition approach requires no semantic understanding or manual annotation, enabling fully automatic discovery of format nesting relationships during live fuzzing campaigns.


\subsubsection{Grammar-guy}

Grammar-guy's goal is to write grammars for any given target program by looking at the harness and deriving a Nautilus-compatible Python grammar from it. The grammar is then used to generate inputs, which are forwarded to Artiphishell's fuzzing components.
To this end, we implemented a set of LLM-based agents that collaborate to perform source code analysis and gain an understanding of its structure and behavior. Based on this understanding, the agents build a grammar that represents the derived input structure and their corresponding data fields. The knowledge accrued in this process is stored and shared between the grammar components using the analysis graph.
We initially developed Grammar-guy on top of Grammarinator \cite{grammarinator}, which uses ANTLR grammars for input generation. ANTLR grammars are however restricted in their ability to represent regular expressions, a limitation the LLM could not handle, leading to the generation of unsupported grammars. Consequentially, we made the decision to switch out input generation to Nautilus \cite{nautilus}, which supports both JSON and Python based grammar representations. After initially working with Nautilus's JSON representation, we adopted the more versatile Python format. This allowed Grammar-guy to generate grammars that contain comprehensive regular expressions and generally proved to be more flexible by also simplifying the expression of byte strings and other data through the possibility of outsourcing terminal data generation to helper functions.
Consider the example initial grammar for the \verb|sqlite_rtree_tables_fuzzer| corresponding initial grammar proposed by Grammar-guy below:
\TBD also put this in appendix?
% Put in appendix aswell?
% \begin{figure*}
%  \begin{minted}[fontsize=\footnotesize, baselinestretch=1.0, frame=single, framesep=3mm]{python}
% # Nautilus grammar for fuzzing the SQLite R*Tree harness with circle query support.
% # This grammar generates input data for the harness, which expects a leading operation selector byte,
% # followed by operation-specific binary data (floats and ints in little-endian format).
% # The four operations are:
% #   0: Insert (id:int, minX:float, maxX:float, minY:float, maxY:float)
% #   1: Select (minX:float, maxX:float, minY:float, maxY:float)
% #   2: Delete (id:int)
% #   3: Select with circle (x:float, y:float, radius:float)
% #
% # The grammar maximizes coverage by generating all valid operation types and edge cases for numeric values.

% # Entry point: one operation per input
% ctx.rule("START", b"{OPERATION}")

% # The first byte selects the operation (0=Insert, 1=Select, 2=Delete, 3=Circle Select)
% ctx.rule("OPERATION", b"{OP_INSERT}")
% ctx.rule("OPERATION", b"{OP_SELECT}")
% ctx.rule("OPERATION", b"{OP_DELETE}")
% ctx.rule("OPERATION", b"{OP_CIRCLE_SELECT}")

% # Insert: 0x00 + id (int32) + minX (float32) + maxX (float32) + minY (float32) + maxY (float32)
% ctx.rule("OP_INSERT", b"\x00{ID}{FLOAT}{FLOAT}{FLOAT}{FLOAT}")

% # Select: 0x01 + minX (float32) + maxX (float32) + minY (float32) + maxY (float32)
% ctx.rule("OP_SELECT", b"\x01{FLOAT}{FLOAT}{FLOAT}{FLOAT}")

% # Delete: 0x02 + id (int32)
% ctx.rule("OP_DELETE", b"\x02{ID}")

% # Circle Select: 0x03 + x (float32) + y (float32) + radius (float32)
% ctx.rule("OP_CIRCLE_SELECT", b"\x03{FLOAT}{FLOAT}{FLOAT}")

% # ID: 4 bytes, signed int32, little-endian
% def id_le_bytes(INTVAL: bytes) -> bytes:
%     # INTVAL is 4 bytes, treat as little-endian int32
%     return INTVAL

% ctx.script("ID", ["INTVAL"], id_le_bytes)

% # INTVAL: generate edge-case and random int32 values
% def int32_edge_cases() -> bytes:
%     import struct
%     # Edge cases: 0, 1, -1, INT_MAX, INT_MIN, random
%     cases = [
%         0,
%         1,
%         -1,
%         0x7fffffff,
%         -0x80000000,
%     ]
%     import os
%     # Add a random int32
%     randval = int.from_bytes(os.urandom(4), "little", signed=True)
%     cases.append(randval)
%     # Pick one at random
%     import random
%     val = random.choice(cases)
%     return struct.pack("<i", val)

% ctx.script("INTVAL", [], int32_edge_cases)

% # FLOAT: 4 bytes, IEEE754 float32, little-endian
% def float32_edge_cases() -> bytes:
%     import struct
%     # Edge cases: 0.0, -0.0, 1.0, -1.0, inf, -inf, nan, FLT_MAX, FLT_MIN, random
%     import math
%     cases = [
%         0.0,
%         -0.0,
%         1.0,
%         -1.0,
%         float("inf"),
%         float("-inf"),
%         float("nan"),
%         3.4028235e+38,   # FLT_MAX
%         -3.4028235e+38,  # -FLT_MAX
%         1.17549435e-38,  # FLT_MIN (smallest positive normal)
%         -1.17549435e-38, # -FLT_MIN
%     ]
%     import os
%     # Add a random float32
%     randval = struct.unpack("<f", os.urandom(4))[0]
%     cases.append(randval)
%     # Pick one at random
%     import random
%     val = random.choice(cases)
%     return struct.pack("<f", val)

% ctx.script("FLOAT", [], float32_edge_cases)
% # BAZINGA
% \end{minted}    
% \end{figure*}
% Cool initial grammar - problem : Harness and grammar are huge. I don't have logs for other initial grammars
%
After the initial grammar is generated, Grammar-guy goes through a fix/refinement phase. The agent passes the proposed grammar to the generator module of Nautilus, and collects a positive or negative format check. We extended the reporting capabilities of Nautilus to provide extra information so that the error description can be used to fix the grammar. Given the generators feedback, Grammar-guy tries to fix each grammar a limited number of times (e.g., 8). In case of failure, Grammar-guy discards the broken grammar and generates a new one from scratch, repeating the same process until a valid initial grammar is generated.

Once the initial generation step yields a working grammar, Grammar-guy generates a set number of seeds (e.g., 200). For each seed, we  then execute the instrumented target to collect coverage information and a report of the accumulated coverage across all seeds, which allows Grammar-guy to determine the reachability of each generated grammar. We store retrieved coverage information in a data structure that associates the grammars with the achieved function coverage for later reuse.

%We keep two data structures:
%grammar_dict: key = func_name value=[grammar1, grammar2]
%grammar_coverage: key = grammar value=[FunctionCoverageMap]

Once a grammar has been successfully derived, coverage was collected and information about its reachability was stored, Grammar-guy enters an improvement loop that finishes after when one of two conditions is reached: a) the budget is exceeded and no new budget is scheduled (in this case, we stall indefinitely, waiting for new budget), b) the campaign is over and ARTIPHISHELL terminates Grammar-guy.
For grammar improvement, Grammar-guy uses three different strategies:
%
\begin{enumerate}
\item \textbf{random}: Grammar-guy selects a function that either (i) has not been reached or (ii) has been covered partially and tries to modify the grammar to reach into that function.
\item \textbf{callable-pair}: Grammar-guy looks for functions that are already reached by the grammar and tries to call functions that are invoked from there. More precisely, Grammar-guy creates a list of pairs (caller, callee) and randomly selects one.
The conditions for reaching the new functions are determined by invoking a specialized agent that inspects both the caller and callee and provides a detailed report on the conditions to reach the callee function.
\item \textbf{extrapolation}: Grammar-guy takes a previously generated grammar and uses the LLM to identify key features encoded in the grammar as well as missing features that are currently not represented in the grammar. This step aims to leverage all of the LLM's internalized knowledge about the target format.
\end{enumerate}
%

The new grammar is again validated and fixed if necessary. Valid grammars are again used to generate more seeds, which are serve to collect coverage and evaluate the reachability of the grammar in the program under test.

When a newly generated grammar yields coverage on functions that were previously unseen (not just new parts of a function), Grammar-guy stores this information in the analysis graph, allowing other components to make use of it.
In addition, the corresponding seeds are passed to the \seedsync agent.

\subsubsection{Grammar-agent}

Grammar-guy's partially randomized, semi-guided algorithm to expand coverage based on static reachability allows it to incrementally progress through the coverage space of the target application.

However, LLMs are often able to bridge large gaps in coverage space based on the semantic meaning of a given function or other internalized knowledge about the program.
For example, if given a harness that parses HTTP requests, LLMs can often directly one-shot grammars that are able to reach most parts of the header and body parsing logic due to their pretrained knowledge of the file format and valid fields. Furthermore, they are also often able to identify interesting functions that are likely reachable via the given harness based on the calling contexts and the names of these functions and knowledge that was learned during the training phase.
As such, grammar-agent represents a component that performs fully LLM-guided exploration of the target application space.

        At a high-level the grammar-agent maintains the following state throughout its execution:
        
        \begin{enumerate}
        \item A set of "goal"-functions that the agent is attempting to reach. These can be added both by the agent itself as well as dynamically based on pipeline inputs.
        For example, in delta-mode tasks we automatically initialize the goal functions list to contain all functions that were modified in the given diff.
        Lastly, if the agent appears to not be making progress for extended periods of time, we will inject goal functions into the goals as well to attempt to divert the agent to discover other code.
        
        \item The full map of all files and functions covered by the agent so far, including the grammars that reached them.
        Any newly discovered files and functions are provided as notifications to the agent to alert it that new coverage has been reached.
        Furthermore, this allows the agent to fully identify the functions that were reached.
        For example, if a function is defined in a C header file and included multiple times, the agent will receive the unique identifier for the exact function that was reached.
        
        \item All grammars that discovered previously unreached code.
        When a grammar reaching new functions is discovered and added to this corpus, it is first uploaded to the analysis graph and linked to the nodes representing all the functions it reached.
        Secondly, these grammars are also automatically distributed across the nodes in our CRS by the \seedsync agent for use by grammar-composer and the grammar-fuzzing instances via the revolver-mutator library.
        
        \item A limited set of LLM-written "memories" that the LLM can use to remember information permanently.
        
        \end{enumerate}
        
        The system prompt to the agent includes general instructions describing the task to write nautilus grammars to maximize the number of functions reached in the target application, as well as formatting instructions and best practices for good grammar writing.
        The user prompt includes the harnesses available to the agent, the current state of the "goal report", as well as the current set of "memories" the LLM stored.
        
        Grammar-agent generally operates by exploring the code with source-reading tool-calls, maintaining its goal list by adding and giving up on goals, and testing its grammars by retrieving their coverage reports.
        
        \begin{itemize}
            \item \textit{add\_goal\_function}: Adds a function by fuzzy-matching the LLM-provided identifier against the available indexed functions. If a unique matching identifier is found, this function is added to the goal functions and immutably added to the "Goal Report".
            \item \textit{give\_up\_on\_goal}: Allows the agent to give up on a given function by providing a reasoning message describing why it believes the function to be unreachable. Generally this happens in case the agent repeatedly fails to reach a function.
            \item \textit{goal\_report}: Retrieves an updated state of the goal report.
            \item \textit{find\_function}: Retrieves the code of a given function via an LLM-provided identifier.
            \item \textit{check\_grammar\_coverage}: The LLM provided nautilus grammar is used to produce N inputs (currently 20), whose coverage is then evaluated. The LLM also provides a list of "functions\_of\_interest" whose coverage reports are included in the response to the LLM. Lastly, if any new functions or files are reached, the corresponding notifications are returned as well.
            \item \textit{remember}: A given memory is added to the list of memories maintained in the user-prompt.
            \item \textit{grep\_sources}: Allows the LLM to provide an expression which is used to grep across the source code. The matching lines and file names are returned.
            \item \textit{get\_file\_content}: Given a file path and a [start,end)-line-range, the corresponding lines of the file are returned.
            \item \textit{get\_functions\_in\_file}: Return all indexed function identifiers in an LLM-provided file.
            \item \textit{get\_function\_calls}: Given a function identifier, return all functions
            \item \textit{get\_files\_in\_directory}: Given a directory path, return all files in the target source directory.
        \end{itemize}

\subsubsection{Grammaroomba}

While the exploration components of our grammar-based fuzzing pipeline serve to improve general reachability, they tend to miss parts of functions that require more intricate grammar constructions. 
To make up for this, we developed Grammaroomba, a local exploration agent that leverages information from the analysis graph to enhance grammars such that coverage in already hit functions is maximized.

Grammaroomba waits for the analysis graph to be populated and then consults the analysis graph to retrieve previously hit functions and grammars associated with them. These pairs additionally come with function coverage information, allowing us to sort out functions that have already been fully covered. Grammaroomba then makes use of a multi factor code complexity ranking based on cyclomatic complexity, nesting depth, exception handling and more. We do so, to allow Grammaroomba to prioritize functions that are assumed to be more prone to security issues and hence have higher relevance for exploration. We select the highest ranked functions to be coverage optimized and explored further.

Grammaroomba leverages the same tools and improvement strategies as Grammar-agent while restricting itself on exploring coverage in one function only. Grammaroomba is also constrained to a fixed amount of coverage runs before skipping over a function and moving on to the next. Other than Grammar-guy and Grammar-agent, Grammaroomba can decide to immediately drop functions without attempting to explore them even though they were selected for improvement by our ranking. The functions that were selected to be optimized are tracked within Grammaroomba to avoid redundant work and new information is retrieved from the analysis graph, ranked and added to the improvement queue in fixed intervals. 

\subsubsection{\revolver}

\revolver is a NAUTILUS-based mutator library exposed to AFL++ and Jazzer. Grammars are synced to fuzzers from other components, e.g., Grammar Guy, and ingested by \revolver to allow the fuzzers to mutate according to those grammars. Inputs to \revolver are serialized in the Rust Objection Notation (RON) format, which contain a grammar string, grammar hash, and a serialized representation of a Nautilus tree representing the structure of the rules in the tree and their values. A syncing component called Watchtower is used to watch for new grammars and produce RONs from the python grammar files, as well as to produce byte representation of new inputs and crashes found by a fuzzer. 
\subsection{\discoveryguy}

The \discoveryguy agent takes as input a \emph{sink function}, which has been identified as potentially containing a bug, tries to identify the actual vulnerability, and, if successful, creates a Python script that generates the corresponding crashing seed. 

\discoveryguy is composed of several sub-agents that cooperate to achieve its overall goal. 
%as shown in Figure~\ref{fig:discoveryguy}.
Note that \discoveryguy is invoked repeatedly on each function that has been identified as interesting by \codeswipe.

%\begin{figure*}[h]
%  \centering
%  \includegraphics[width=\linewidth]%{figures/discoveryguy.png}
%  \caption[\discoveryguy]{The \discoveryguy algorithm, which shows the control flow among the various sub-agents.}
%  \label{fig:discoveryguy}
%\end{figure*}

\discoveryguy is also invoked whenever a SARIF report is received.
In this case, in addition to the sink function(s), \discoveryguy receives as input the type of vulnerability identified by the static analysis tools.
At the end of the process, if a crashing seed is found, the SARIF report is considered valid.

Finally, when a patch is generated, the patched function is passed to \discoveryguy together with the original code, and the agent is asked to try to bypass the patch and still exploit the vulnerability.
If \discoveryguy can find a crashing seed, the patch is invalidated.

These three use cases follow a similar workflow, which is described hereinafter.

As the first step, \discoveryguy analyzes the call graph contained in the \analysisgraph to identify the harnesses that could be used to reach the sink function.
This analysis includes a graph reduction step that simplifies the structure of the call graph to make its size consumable by the LLM (i.e., sets of nodes with the same pre- and post-dominator nodes are merged).
If no harnesses are found, possibly because of the incompleteness of the information extracted from the target's code, \discoveryguy asks the LLM to produce an educated guess of which harnesses among those available are most likely to reach the sink function. 

Then, the information about the suitable harnesses is passed to \discoveryguy's \jimmypwn agent, which is based on \claudefour.
\jimmypwn is tasked with identifying the vulnerability and producing a program that generates a crashing seed. \jimmypwn has access to a grepping tool that can find specific patterns in the overall code base.

If successful, \jimmypwn generates a report that contains the following information: (i) the harness that is used to reach the sink function and the associated call trace, (ii) the relevant path conditions that are necessary to drive the harness input to the location of the vulnerability, and (iii) a script that, when executed, generates the crashing seed.

This report is passed to the \seedgeneration agent (based on \ofourmini), which refines the original seed generation script to verify that it actually crashes the function.
This is necessary because sometimes, even if the vulnerability is successfully identified, the script that generates the crashing seed is not effective.
For example, the \jimmypwn agent might have identified an overflow that affects a 256-character buffer, but because of the LLM imprecision, the seed-generation script creates an input that is only 128-byte long.
Therefore, \seedgeneration iterates a few times on the seed generation script to improve its effectiveness.

If this process does not converge to a working script, then the information about the failed script (why and where it fails) is sent back to the \jimmypwn agent requesting the generation of a new report for this vulnerability.
This feedback loop is repeated a few times until an effective seed generation program is created (and in case of failure, the whole process is aborted).
Note that in the case of a SARIF report analysis, this process is repeated for additional rounds, as there is a higher confidence that the sink function contains a vulnerability. 

% Understand the role of the crash checker


\subsection{LLuMinar}
\label{lluminar}
The LLuMinar agent scans all functions and methods reachable from any harness entry point, using our customized reasoning model to autonomously retrieve relevant context and determine whether a function is vulnerable. If a function is predicted to be vulnerable, the model outputs the predicted CWE type along with a natural language explanation. If the function is determined to be benign, the model outputs reasoning that explains why it poses no security risk.

\subsubsection{Agent scaffold}
\label{sec:lluminar_agent}

LLuMinar is designed as a lightweight vulnerability scanning agent that rapidly identifies suspicious functions for downstream heavy-weight components (such as DiscoveryGuy) to target for PoC generation. 
Given the time constraints of the AIxCC competition, we developed a light-weight agent scaffold that incorporates one function retrieval tool with a maximum of five tool invocations to enable fast scanning.

We define target functions as all functions or methods reachable from any harness entry point, meaning functions for which there exists a call graph path from the harness entry point defined in ~\cref{sec:oss-fuzz} in the Analysis Graph database. 

For each target function and its corresponding reachable harnesses, we extract contextual information by sampling three random paths in the call graph from each harness entry point to the target function and providing the implementation of each intermediate function along these paths as initial context. 
The model is equipped with a tool that can retrieve the implementation of any function given its name, enabling the model to fetch any missing context autonomously. 
Once the model determines that sufficient information has been collected, it outputs its vulnerability prediction along with reasoning. 

\subsubsection{Model Fine-tuning}

For fine-tuning, we use Qwen2.5-7B-Instruct~\cite{yang2024qwen2} as our base model and apply standard supervised fine-tuning on training data constructed from the ARVO~\cite{mei2024arvo} dataset through four carefully designed data sources:

\textbf{(1) Vulnerable and patched function pairs:} For each instance in ARVO, we extract both the pre-patch (vulnerable) and post-patch (benign) versions of the same function with their corresponding stack traces as context, collecting reasoning from o3. This paired data enables our model to learn the subtle differences between vulnerable and secure implementations of the same functionality, developing fine-grained reasoning capabilities for vulnerability detection.

\textbf{(2) Benign functions:} We randomly select benign target functions and their call graph paths as context, with reasoning also distilled from o3. Since the paired vulnerable/patched data primarily focuses on functions that have experienced security issues, we supplement our training with inherently safe functions to help the model learn the typical properties of non-vulnerable code.

\textbf{(3) Proof-of-concept generation data:} For each vulnerability, we provide prompts including the stack trace, ground truth PoC, and patch, with reasoning distilled from Claude-4. We believe that training the model on proof-of-concept (PoC) generation tasks enhances its vulnerability detection performance. The rationale is that PoC generation inherently requires the model to reason concretely about how inputs propagate along execution paths and precisely identify the conditions triggering vulnerabilities. By including PoC generation data in training, we encourage our model to adopt this concrete, simulation-based reasoning style rather than relying merely on superficial indicators of code vulnerability, such as code smells or simplistic patterns.

\textbf{(4) Agent-based interaction data:} We collect trajectories  using the agent described in ~\cref{sec:lluminar_agent} with o3. This data source specifically enhances our model's ability to effectively utilize our retrieval tool and make informed decisions about when and what additional context to gather.

For all data sources, we applied rejection sampling to retain only samples with correct results. To further increase the proportion of high-quality training data, we adopted a best-of-N strategy: for each data point, we sampled multiple model responses to improve the likelihood of obtaining correct answers. When multiple correct responses were available, we selected the one with the most concise reasoning to promote efficiency in the training data.

\subsection{QuickSeed}
\label{quickseed}
QuickSeed is an LLM fuzzing seed generation tool for Java projects.
The key insight for discovering vulnerabilities in the Java projects is that 
Jazzer's sanitizers hook up specific Java API calls to detect the crash.
For example, to detect OS command injection, Jazzer monitors the invocation of \textit{java.lang.ProcessBuilder.start}.
This means that if we find all the API calls used by the sanitizers  in the project, we have all the possible crash locations.
We can use these code locations as sinks and analyze how they can be reached from the harnesses.
In other words, we can construct the call graph, find all paths that lead to these sinks.
Besides the APIs used by the sanitizers, we also analyze Java CVEs and find more promising sinks.

After getting all paths to all sinks, QuickSeed takes one sink and one path that can reach the sink as the input and asks LLM to generate a seed that can crash the sink.
This is the seed generation agent.
At this stage, we only give LLM all the function code on the call path, the harness code, and the sink code.
After LLM generates a seed, we run the generated seed against the target and see whether it crashes.
If it crashes, we go on to the next path.
Otherwise, we collect the execution trace of the run and see how many functions on the path the seed has covered.
We then invoke a feedback LLM agent with the coverage information and let the LLM figure out a way to bypass the potential blocker that might exist in the stuck function.
We give the source code retrieval ability to the feedback agent so that the agent can explore the code base.
The reason we do this is that, in many cases, the blocker is not on the call path.
By giving the LLM agent the freedom to explore the code base, we can avoid the complicated data flow analysis and rely on the LLM's ability to analyze the code and find the points of interest critical to reach the sink function.

The biggest challenge is the false negative of the call analysis. 
For example, the static analysis cannot resolve a reflection call.
Our solution to resolve this is to run a warm-up LLM agent and generate some seeds for every harness.
Then we trace these seeds and construct a dynamic call graph.
And also for every seed that is traced, we use the tracing result to complete the call graph.
Another challenge is that we can have too many paths.
To mitigate this issue, we utilize the ranking of CodeSwipe to prioritize the high-value sinks. 
And also, if for one sink there are many paths that can reach it, we will query the dynamic call graph we have and prioritize the path in the dynamic call graph over the static call graph.
\subsection{\aijon}
\label{aijon}

\aijon is a fuzzer that implements IJON\cite{aschermann2020ijon} annotations on top of AFL++.
It contains an LLM agent which inserts these annotations into the target's source.

\aijon takes as input a list of functions that are identified by \codeswipe to possibly contain a vulnerability.
It then uses the LLM agent to analyze the source code of this single function, so as to reduce the LLM context, and inserts IJON style annotations into its source code.

The LLM is instructed to strategically select annotations that can guide the fuzzer towards the vulnerability specified in the \codeswipe report.
These annotations provide valuable feedback to AFL++ that can guide it very quickly to the target locations and find a crashing input if any for the specified vulnerability.

However, using such annotations can lead to collisions within the bitmap of AFL++.
In order to address this problem, we optimize the instrumentation step of AFL++ by limiting the instrumentation to a select list of functions.
\aijon uses the callgraph contained in the \analysisgraph to identify the functions on the path from the harness' entry point to the target function.
It then instructs AFL++ to only instrument the functions that are present in these paths.
This allows AFL++ to focus on the vulnerability in the target function without exploring other sections of the target program and reduces the possibility of collisions in the bitmap.

\aijon also uses the \analysisgraph to generate a starting seed corpus for fuzzing.
Other fuzzers, such as \afluzzer, could have accidentally generated inputs that reach some of the functions in the call path to the target function.
These inputs are optimal candidates for being used as seeds, since they have triggered a function close to the target function in the callgraph.
\aijon collects such inputs from \analysisgraph and uses them as the starting seed corpus.

Any interesting inputs or crashes found by \aijon are synced with \afluzzer.
This allows \afluzzer to replicate the crashes on the target without annotations and acts as a verification step for the crashes detected by \aijon.

\subsection{\povguy}

\povguy is responsible for taking a crashing seed and a harness and run the seed in a setup that is as close as the canonical target as possible.

First, \povguy pulls the build with the relevant santizer from \builder and runs the input five time to ensure that the input trggers the sanitizer reliably (this is necessary to avoiding that the evaluators fail to reproduce the attack and give a negative score).

The \povguy parses the crashes and determines all the aspects for crash, such as the fact that it was caused by a read or by a write operation.
It also parses all the stack traces for a crash and extracts the various crash-related information (note that a sanitizer can generate multiple stack traces).

Then, \povguy extracts the information necessary for the deduplication of the crash. 
\povguy has four different types of deduplication approaches:
\begin{enumerate}
\item Dedup OSS-Fuzz: it is the deduplication token used by libfuzzer, which uses the names of the last three functions in stack;
\item Dedup token full: concatenates everything in a stack traces;
\item Dedup token Shellphish: similar to the OSS-Fuzz one, but it only uses the functions that are indexed in the source code; this avoids adding to the token library functions such as memcopy, reducing the noise;
- Dedup token organizers: it uses only the stack of the crashing state associated with the triggering of the sanitizer and generate a token. In the case of Jave, this is an ``instrumentation key,'' which is the sorted list of all classes where at least a function was executed. 
\end{enumerate}
\povguy uploads this information to the \analysisgraph.

Then, \povguy reads all the historical dedup tokens from the \analysisgraph  and determines if the current dedup token is a duplicate of an existing one.
%For this task, \povguy uses the method from the organizers which uses some edit-distance based similarity.
If there is a similarity, a dupliocation edge is added between the tokens.

In delta mode, \povguy makes sure that the input crashes in the application with the delta and not in the base application, ensuring that the triggered vulnerability is in the added code.




If the crash is not a duplicate it goes to the submitter agent which submitted to the scoring system'

